These notes introduce probabilistic methods as a way of reasoning about uncertainty. The course begins with sample spaces and events, then develops counting techniques, probability distributions, standard distribution functions, conditional probability, and finally descriptive statistics based on generated data.

The central aim is to connect formal probability with interpretation. A probability calculation should not be treated as a mechanical substitution into a formula. It should begin with a clear description of the random experiment, the elementary outcomes, the events or random variables of interest, and the assumptions that make a chosen model appropriate.
